% ============================================================ Capítulo 10
\chapter{Instalación y configuración de OpenCode}\label{cap:opencode}

\section{Qué es y cómo se usa en el proyecto}
OpenCode es una consola agéntica de código abierto, desacoplada del proveedor de modelos y extensible mediante
MCP~\cite{opencode}. Incluye dos agentes primarios que se alternan con la tecla \texttt{Tab}: \textbf{plan}
(solo lectura, planifica) y \textbf{build} (ejecuta cambios). Cada acción de herramienta se rige por
\texttt{allow}, \texttt{ask} o \texttt{deny} (Clase~3). Los \emph{forks} previos (\texttt{sst/opencode},
\texttt{opencode-ai/opencode}) están archivados; el repositorio activo es \texttt{anomalyco/opencode}.

\section{Instalación y verificación}
\begin{lstlisting}[style=ps]
npm install -g opencode-ai          # Windows (alternativas: scoop install opencode, choco install opencode)
opencode --version
cd C:\fabrica-nacional-autopartes
opencode                            # abre la interfaz en el repositorio; /init analiza el proyecto
\end{lstlisting}

\section{Proveedor LLM y cuenta de Claude}\label{sec:proveedor-opencode}
\begin{nota}[Suscripción Claude y OpenCode]
Una suscripción Claude Pro o Max \textbf{no} incluye créditos de API. Según los términos de Anthropic
actualizados en febrero de 2026, los tokens OAuth de las cuentas Free, Pro y Max solo pueden usarse en Claude
Code y en claude.ai; OpenCode retiró el inicio de sesión con suscripción de Anthropic tras requerimientos
legales~\cite{register-anthropic}. Existen complementos comunitarios que lo reintroducen; su uso contraviene esos
términos y no se recomienda. Opciones válidas para OpenCode: Ollama local, una clave de API con facturación por
uso (Consola de Anthropic u otro proveedor) o un servicio compatible. Verificar en \cmd{opencode auth list} y
\texttt{/models} qué proveedores quedan disponibles en la versión instalada.
\end{nota}

\begin{lstlisting}[style=ps]
opencode auth login                 # elegir proveedor e ingresar la clave (se guarda fuera del repositorio)
opencode auth list
ollama pull qwen2.5-coder:14b       # modelo local para tareas de código (>= 14B recomendado)
\end{lstlisting}
Los modelos locales pequeños como \texttt{llama3.2:3b} sirven para la aplicación, pero no para el bucle
agéntico de edición de código: la documentación de Ollama recomienda una ventana de al menos 64k tokens para
OpenCode~\cite{ollama-opencode}.

\section{Configuración del proyecto}
\ruta{opencode.json} declara las instrucciones (\texttt{AGENTS.md}), el proveedor Ollama, los permisos, cuatro
subagentes y los servidores MCP:
\lstinputlisting[style=json,firstline=1,lastline=31]{../../opencode.json}
\begin{table}[H]\centering\small
\caption{Agentes especializados del proyecto}\label{tab:agentes}
\begin{tabularx}{\linewidth}{@{}l X l@{}}\toprule
Agente & Responsabilidad & Herramientas\\\midrule
analista-requisitos & Matriz de trazabilidad, especificaciones OpenSpec & sin bash; edición con aprobación\\
ingeniero-datos & Scraping y pipeline Medallion & lectura, edición, pruebas\\
ingeniero-llm & Prompts, cadenas, salvaguardas & lectura, edición, pruebas\\
revisor-calidad & Pruebas, lint, tipos, capas; reporta sin editar & sin \texttt{write}/\texttt{edit}\\\bottomrule
\end{tabularx}
\end{table}

\section{Contexto, reglas, ejecución y revisión}
El contexto del proyecto se carga desde \texttt{AGENTS.md} y \ruta{openspec/config.yaml}. Las reglas exigen
ejecutar pruebas, lint y tipos antes de declarar una tarea terminada. Flujo recomendado: (1) agente
\textbf{plan} con la especificación; (2) revisión humana del plan; (3) agente \textbf{build}; (4)
\cmd{git diff} y revisión; (5) subagente \texttt{@revisor-calidad}; (6) \emph{commit}. La depuración se apoya
en la salida de pytest pegada en la conversación y en la regla ``corregir la causa, no la prueba''.

% ============================================================ Capítulo 11
\chapter{Instalación y configuración de Claude Code}\label{cap:claudecode}

\section{Instalación y autenticación}
\begin{lstlisting}[style=ps]
irm https://claude.ai/install.ps1 | iex     # instalador nativo (sin administrador)
# alternativa: winget install Anthropic.ClaudeCode
claude --version
cd C:\fabrica-nacional-autopartes
claude                                      # primer inicio: autenticación en el navegador
claude doctor
\end{lstlisting}
Claude Code funciona de forma nativa en Windows~10 (1809+) y Windows~11; Git for Windows es recomendable porque
aporta Git Bash para la herramienta de shell~\cite{claude-code}.

\section{Suscripción frente a acceso por API}
\begin{table}[H]\centering\small
\caption{Métodos de autenticación}\label{tab:auth}
\begin{tabularx}{\linewidth}{@{}l X X@{}}\toprule
Método & Facturación & Alcance\\\midrule
Suscripción Pro/Max/Team/Enterprise & Cuota del plan con límites de uso & Claude Code y claude.ai\\
Consola de Anthropic (clave de API) & Pago por token & Claude Code, SDK, aplicaciones propias y herramientas de terceros\\
Proveedores en la nube (Bedrock, Vertex) & Facturación del proveedor & Entornos corporativos\\\bottomrule
\end{tabularx}
\end{table}
La aplicación del proyecto (proveedor \texttt{anthropic}) requiere una clave de API en
\texttt{FNA\_ANTHROPIC\_API\_KEY}; la suscripción no la sustituye.

\section{Configuración del repositorio}
\begin{itemize}
  \item \textbf{CLAUDE.md}: importa \texttt{@AGENTS.md} (fuente única de reglas) y añade lo específico de Claude Code.
  \item \textbf{Subagentes} (\ruta{.claude/agents/*.md}): nueve roles con herramientas restringidas y el modelo
        fijado en el campo \texttt{model} del \emph{frontmatter} (sección~\ref{sec:modelos-especialidad}); espejo
        en \texttt{opencode.json}.
  \item \textbf{Skills} (\ruta{.claude/skills/}): copia sincronizada de \ruta{skills/}.
  \item \textbf{Hooks}: tras \texttt{Edit|Write|MultiEdit} se ejecuta \cmd{uv run ruff format}.
  \item \textbf{MCP}: \ruta{.mcp.json} registra \texttt{fna-conocimiento}; \texttt{enabledMcpjsonServers} limita
        los servidores habilitados.
\end{itemize}
\lstinputlisting[style=json,firstline=1,lastline=41]{../../.claude/settings.json}

\section{Modelos por especialidad}\label{sec:modelos-especialidad}
Claude Code CLI es la consola principal. Cada subagente declara \texttt{model: opus}, \texttt{sonnet} o
\texttt{haiku}; si el campo se omite, hereda el modelo de la sesión. La variable
\texttt{CLAUDE\_CODE\_SUBAGENT\_MODEL}, cuando se define, reemplaza el campo de todos los subagentes, por lo que
debe quedar sin definir para respetar la asignación.
\begin{table}[H]\centering\small
\caption{Asignación de modelos de Claude por especialidad}\label{tab:modelos-especialidad}
\begin{tabularx}{\linewidth}{@{}P{3.4cm} X P{1.4cm}@{}}\toprule
Subagente & Especialidad & Modelo\\\midrule
analista-requisitos & Requisitos, trazabilidad, OpenSpec & Opus\\
arquitecto-sistema & Arquitectura, ADR, integración API/MCP/PostgreSQL/visor & Opus\\
arquitecto-mecanico-3d & Modelo de componentes, articulaciones, límites, restricciones, física simplificada & Opus\\
revisor-calidad & Revisión independiente sin permisos de edición & Opus\\
ingeniero-datos & Scraping, Medallion, métricas y gráficas & Sonnet\\
ingeniero-llm & Prompts, LangChain, RAG del simulador & Sonnet\\
ingeniero-3d-frontend & Generación GLB, API del simulador, visor 3D, animaciones, controles, pruebas & Sonnet\\
documentador & README, guías, manual LaTeX & Sonnet\\
verificador-rapido & Verificaciones repetitivas (aprobado/reprobado) & Haiku\\\bottomrule
\end{tabularx}
\end{table}
\begin{lstlisting}[style=txt]
---
name: arquitecto-mecanico-3d
description: Diseña el modelo de componentes del simulador 3D ...
tools: Read, Grep, Glob, Edit, Bash
model: opus
skills: fna-simulador-3d
---
\end{lstlisting}
La sesión principal se inicia con \cmd{claude --model opus} para planificar y cambia a \cmd{/model sonnet} para
implementar. El alias combinado de planificación e implementación se confirma con \cmd{/model} en la versión
instalada antes de usarlo. En OpenCode, los mismos roles usan identificadores de la API
(\texttt{anthropic/claude-opus-5-5}, \texttt{anthropic/claude-sonnet-5-5}, \texttt{anthropic/claude-haiku-5-5}),
que requieren clave de API; se verifican con \cmd{opencode models}.

\section{Comandos y validación de cambios}
\begin{lstlisting}[style=ps]
claude                       # sesión interactiva
/agents                      # lista y edita subagentes
/mcp                         # estado de servidores MCP
claude mcp list
claude -p "Ejecutar la skill fna-pruebas y resumir el resultado"   # modo no interactivo
\end{lstlisting}
Todo cambio se valida con la batería de la sección~\ref{sec:bateria} y con el subagente
\texttt{revisor-calidad}, que no tiene permisos de edición.

% ============================================================ Capítulo 12
\chapter{Integración de modelos LLM}\label{cap:llm}

\section{Selección del modelo}
La Clase~2 compara modelos propietarios (arranque inmediato, costo por token, datos en el proveedor) y de pesos
abiertos (costo fijo de hardware, datos locales, auditables). Para el Módulo~1 se elige \textbf{Ollama con
\texttt{llama3.2:3b}} por defecto: ejemplo explícito del taller, costo cero por token y privacidad. El
proveedor se cambia sin modificar código:
\begin{lstlisting}[style=ps]
uv run fabrica ask "¿Qué productos ofrece Fanalvías?" --proveedor ollama --modelo llama3.2:3b
uv run fabrica ask "¿Qué productos ofrece Fanalvías?" --proveedor openai --modelo gpt-4o-mini
uv sync --extra anthropic
uv run fabrica ask "¿Qué productos ofrece Fanalvías?" --proveedor anthropic --modelo <modelo-disponible>
uv run fabrica ask "¿Qué productos ofrece Fanalvías?" --proveedor extractivo   # línea base sin LLM
\end{lstlisting}

\section{Configuración crítica: ventana de contexto}
El taller exige consolidar toda la base en el prompt de sistema. Ollama aplica una ventana por defecto
inferior a la capacidad del modelo, por lo que la fábrica de modelos fija \texttt{num\_ctx}
(\texttt{FNA\_OLLAMA\_NUM\_CTX}, 16\,384 por defecto). Si la base supera
\texttt{FNA\_PRESUPUESTO\_TOKENS\_CONTEXTO}, el modo \texttt{auto} cambia a selección BM25
(Figura~\ref{fig:contexto-modos}). Con la semilla, la base completa ocupa unos 3\,700 tokens aproximados.

\begin{figure}[H]\centering
\begin{tikzpicture}[node distance=0.5cm and 0.9cm]
  \node[caja=2.6cm] (q) {Pregunta};
  \node[diamond,draw=fnagris,aspect=2.2,font=\sffamily\scriptsize,fill=oro!10,right=of q] (d) {¿Base $\leq$ presupuesto?};
  \node[clave=3.3cm, above right=0.2cm and 1.1cm of d] (c) {Modo \textbf{completo}: todos los fragmentos (Módulo 1)};
  \node[clave=3.3cm, below right=0.2cm and 1.1cm of d] (s) {Modo \textbf{selección}: k fragmentos BM25 (extensión)};
  \node[caja=2.6cm, right=4.9cm of d] (p) {Prompt con \texttt{<contexto>}};
  \draw[flecha] (q) -- (d); \draw[flecha] (d) |- node[etq,pos=0.25,left]{sí} (c); \draw[flecha] (d) |- node[etq,pos=0.25,left]{no} (s);
  \draw[flecha] (c) -| (p); \draw[flecha] (s) -| (p);
\end{tikzpicture}
\caption{Modos de armado del contexto}\label{fig:contexto-modos}
\end{figure}

\section{Prompt Engineering}
\begin{table}[H]\centering\small
\caption{Versiones de prompt por tarea}\label{tab:prompts}
\begin{tabularx}{\linewidth}{@{}l X X@{}}\toprule
Versión & Técnica & Hipótesis\\\midrule
v1 & Zero-shot básico (línea base) & Sin reglas, el modelo completa con conocimiento propio.\\
v2 & Zero-shot + rol + reglas + delimitación XML + defensa ante inyección + frase de rechazo fija & Las reglas
explícitas reducen alucinaciones y neutralizan órdenes incrustadas.\\
v3 & Salida JSON + \emph{few-shot} (2 ejemplos) + citas verificables + autoevaluación de confianza (Q\&A);
cadena de verificación interna (Resumen) & La estructura permite validar automáticamente la fundamentación.\\\bottomrule
\end{tabularx}
\end{table}
El prompt v3 de Q\&A completo:
\lstinputlisting[style=yaml,firstline=50,lastline=85]{../../src/fabrica_autopartes/llm/prompts/qa.yaml}

\section{Salvaguardas posteriores a la generación}
\begin{enumerate}
  \item \textbf{Bloqueo de inyección}: patrones como ``ignora todas las instrucciones'' detienen la consulta
        antes de invocar el modelo.
  \item \textbf{Análisis JSON tolerante}: extrae el objeto aunque el modelo agregue texto o bloques de código.
  \item \textbf{Verificación de citas}: una cita a un fragmento que no se entregó invalida la fundamentación.
  \item \textbf{Fundamentación léxica}: al menos el 60\,\% de los términos de la respuesta debe aparecer en los
        fragmentos citados.
  \item \textbf{Normalización del rechazo}: cualquier variante de ``no se encuentra en la base'' se convierte en la
        frase oficial con confianza \texttt{ninguna}.
\end{enumerate}

\section{Experimentación documentada}\label{sec:h01}
La bitácora \ruta{docs/prompts/bitacora_experimentos.md} define el protocolo (mismo banco, misma base,
temperatura 0,1, una variable por experimento). La línea base extractiva produjo el hallazgo H-01: el ejemplo
\emph{few-shot} de v3 usaba la misma etiqueta \texttt{<fragmento>} que el contexto real y fue citado como un dato
de la empresa (``La planta principal está en la ciudad X''). Corrección: formato ilustrativo distinto para los
ejemplos y lectura exclusiva del bloque \texttt{<contexto>}; dos pruebas de regresión lo protegen. La lección
aplica a cualquier LLM: los ejemplos deben distinguirse visualmente del contexto.

% ============================================================ Capítulo 13
\chapter{Arquitectura y configuración MCP}\label{cap:mcp}

\section{Conceptos}
El \emph{Model Context Protocol} (MCP) desacopla el modelo de sus capacidades externas. Un \textbf{host}
(OpenCode o Claude Code) orquesta la sesión; por cada servidor crea un \textbf{cliente} que mantiene una
conexión 1:1; el \textbf{servidor} expone tres primitivas: \textbf{Tools} (funciones ejecutables con esquema de
entrada), \textbf{Resources} (datos de lectura identificados por URI) y \textbf{Prompts} (plantillas
reutilizables). Los mensajes siguen \textbf{JSON-RPC 2.0}.

\begin{figure}[H]\centering
\begin{tikzpicture}[node distance=0.5cm and 1.1cm]
  \node[grupo=fnaazul, minimum width=4.6cm, minimum height=3.9cm] (host) {};
  \node[font=\sffamily\footnotesize\bfseries, anchor=north] at (host.north) {Host: OpenCode / Claude Code};
  \node[caja=3.4cm] (llm) at ($(host.center)+(0,0.6)$) {Modelo + bucle agéntico};
  \node[clave=1.6cm, below left=0.5cm and -1.6cm of llm] (c1) {Cliente 1};
  \node[clave=1.6cm, below right=0.5cm and -1.6cm of llm] (c2) {Cliente 2};
  \node[caja=4.2cm, right=2.4cm of host.north east, anchor=north west, minimum height=2.0cm] (s1)
    {\textbf{fna-conocimiento} (propio)\\Tools: buscar\_fragmentos, obtener\_fragmento, estadisticas\_base, validar\_corpus, responder\_pregunta*\\Resources: fna://base/consolidada, fna://base/fuentes\\Prompts: plantilla\_qa};
  \node[externo=4.2cm, below=0.6cm of s1] (s2) {\textbf{Engram} (existente)\\mem\_save, mem\_search, mem\_context\ldots};
  \node[almacen, right=0.7cm of s1] (g) {Gold\\(lectura)};
  \node[almacen, right=0.7cm of s2] (e) {SQLite\\FTS5};
  \draw[flecha] (llm) -- (c1); \draw[flecha] (llm) -- (c2);
  \draw[flecha] (c1.east) to[out=0,in=180] node[etq,above,pos=0.55]{stdio · JSON-RPC} (s1.west);
  \draw[flecha] (c2.east) to[out=0,in=180] node[etq,below,pos=0.55]{stdio} (s2.west);
  \draw[flecha] (s1) -- (g); \draw[flecha] (s2) -- (e);
\end{tikzpicture}
\caption{Arquitectura MCP del proyecto (* deshabilitada por defecto)}\label{fig:mcp}
\end{figure}

\section{Transportes y ciclo de vida}
\textbf{stdio}: el host lanza el servidor como subproceso; stdout transporta exclusivamente JSON-RPC (por eso la
auditoría del servidor escribe en stderr). \textbf{Streamable HTTP}: el servidor escucha en un puerto
(\cmd{uv run fabrica-mcp --http} en \texttt{127.0.0.1:8765/mcp}); admite autenticación OAuth para despliegues
remotos. Ciclo de vida: \texttt{initialize} (negociación de versión y capacidades) $\rightarrow$
\texttt{notifications/initialized} $\rightarrow$ \texttt{tools/list}, \texttt{tools/call},
\texttt{resources/read}\ldots{} $\rightarrow$ cierre de la conexión. La prueba \texttt{test\_sesion\_stdio\_real}
ejecuta este ciclo completo contra el servidor real.

\section{Implementación del servidor propio}
\lstinputlisting[style=py,firstline=68,lastline=98]{../../src/fabrica_autopartes/mcp_server/server.py}

\section{Autenticación, secretos, errores y auditoría}
El servidor local por stdio hereda la identidad del usuario y no requiere credenciales; no lee
\texttt{.env} salvo la configuración de la aplicación. Los errores de dominio (fragmento inexistente) se devuelven
como resultados de error MCP sin detener el servidor. Cada llamada se registra con fecha, herramienta y
argumentos.

\section{Seguridad frente a inyección y aislamiento}\label{sec:permisos}
\begin{itemize}
  \item Herramientas anotadas como solo lectura, no destructivas, idempotentes y sin mundo abierto.
  \item Límites de entrada: consulta $\leq$ 300 caracteres, $k \leq 10$, pregunta $\leq$ 500 caracteres.
  \item El texto del corpus se declara como \emph{datos} en las instrucciones del servidor; un fragmento con
        órdenes incrustadas (inyección indirecta) no obtiene capacidades adicionales porque el servidor no
        expone escritura ni ejecución.
  \item Permisos de lectura, escritura y ejecución separados: lectura MCP permitida sin confirmación en Claude
        Code; edición con confirmación; \texttt{git push}, \texttt{rm -rf} y \texttt{curl} denegados.
\end{itemize}

\section{Registro y prueba funcional}
\begin{lstlisting}[style=ps]
uv run fabrica build
claude mcp list                                    # fna-conocimiento: connected
# Registro manual alternativo (si no se usa .mcp.json):
claude mcp add fna-conocimiento -- uv run fabrica-mcp
uv run pytest tests/integration/test_mcp.py -q     # 6 pruebas, incluida la sesión stdio real
\end{lstlisting}
Ejemplo de invocación desde el agente: ``Usar la herramienta \texttt{buscar\_fragmentos} con la consulta
\emph{garantía} y citar los identificadores devueltos''.

% ============================================================ Capítulo 14
\chapter{Skills y herramientas de especificación}\label{cap:skills}

\section{Verificación de herramientas}
Antes de instalar cada herramienta se verificaron nombre, repositorio, mantenedor, propósito, instalación,
licencia, mantenimiento y riesgos (Tabla~\ref{tab:verificacion}).

\begin{longtable}{@{}P{2.0cm}P{3.6cm}P{4.6cm}P{4.9cm}@{}}
\caption{Verificación de herramientas de gobierno del desarrollo}\label{tab:verificacion}\\\toprule
Herramienta & Repositorio / licencia & Propósito e instalación & Estado y riesgos\\\midrule\endfirsthead
\toprule Herramienta & Repositorio / licencia & Propósito e instalación & Estado y riesgos\\\midrule\endhead
Skills & Formato \texttt{SKILL.md} (frontmatter \texttt{name}, \texttt{description}) & Conocimiento procedimental
cargado bajo demanda; registro \texttt{skills.sh}: \texttt{npx skills add owner/repo} & Portables entre OpenCode,
Claude Code y Cursor. Riesgo: instrucciones o scripts maliciosos en skills de terceros.\\
Engram & \texttt{Gentleman-Programming/engram}, MIT & Memoria persistente de decisiones (Go, SQLite + FTS5) vía
MCP; \texttt{engram setup opencode}; plugin para Claude Code & Activo en 2026. Escribe en
\texttt{\textasciitilde/.engram}; no guardar secretos en memorias.\\
Archify & \textbf{Ambigüedad}: \texttt{tt-a1i/archify} (skill de diagramas) y \texttt{Aryan1718/Archify}
(\texttt{npx archify-cli init}, análisis de repositorio) & Ambas generan documentación arquitectónica. Se
selecciona \texttt{tt-a1i/archify} (\texttt{npx skills add tt-a1i/archify -g}) por funcionar como skill en
OpenCode y Claude Code & Instalación opcional; escanear con SkillSpector antes de usar. Requiere Node.js.\\
OpenSpec & \texttt{Fission-AI/OpenSpec}, MIT & Desarrollo dirigido por especificaciones;
\texttt{npm install -g @fission-ai/openspec@latest}; requiere Node $\geq$ 20.19 & Verificado: versión 1.16.0 validó
las especificaciones del proyecto.\\
SkillSpector & \texttt{NVIDIA/SkillSpector}, Apache-2.0 & Escáner de seguridad de skills (análisis estático y
semántico, salida SARIF); \texttt{uv tool install git+https://github.com/NVIDIA/skillspector.git} & Activo. Regla
del curso: puntuación 0--20 SAFE, 21--50 CAUTION, 51--100 NO INSTALAR.\\
GGA & \texttt{Gentleman-Programming/gentleman-guardian-angel}, MIT & Revisión con IA en \emph{pre-commit} contra
\texttt{AGENTS.md}; \texttt{gga init}, \texttt{gga install} & Implementado en Bash: en Windows requiere Git Bash o
WSL. Envía el diff al proveedor configurado.\\
Medallion & Patrón de arquitectura de datos (no es un paquete) & Bronze/Silver/Gold & Aplicado en
\ruta{pipelines/}; no requiere instalación.\\\bottomrule
\end{longtable}

\section{OpenSpec en el proyecto}
\ruta{openspec/config.yaml} inyecta el contexto del proyecto; \ruta{openspec/specs/} contiene seis capacidades
(extracción de datos, base de conocimiento, tareas LLM, interfaz de prueba, simulador 3D y biblioteca de máquinas) con requisitos en formato
\emph{Requirement/Scenario} y cláusulas WHEN/THEN. \ruta{openspec/changes/} contiene un cambio archivado
(selección de contexto BM25, implementado) y uno propuesto (RAG con embeddings, Módulo~2).
\begin{evidencia}[openspec validate --all (OpenSpec 1.16.0)]
\lstinputlisting[style=txt]{../evidencias/openspec_validate.txt}
La primera validación reportó dos requisitos sin escenario; se corrigieron antes de esta ejecución.
\end{evidencia}

\section{Skills del proyecto}
\begin{longtable}{@{}P{3.4cm}P{5.4cm}P{6.3cm}@{}}
\caption{Skills del proyecto (\texttt{skills/\textless nombre\textgreater/SKILL.md})}\label{tab:skills}\\\toprule
Skill & Responsabilidad & Validación\\\midrule\endfirsthead
\toprule Skill & Responsabilidad & Validación\\\midrule\endhead
fna-requisitos & Analizar requisitos y mantener la trazabilidad & Cada RF con prueba existente\\
fna-codigo & Generar código por capas, tipado y probado & pytest, Ruff, mypy sin errores\\
fna-arquitectura & Revisar capas y ADR & \texttt{verificar\_capas.py} = 0\\
fna-scraping-medallion & Scraping responsable y Medallion & Pruebas de rastreo y \texttt{fabrica validate}\\
fna-validacion-datos & Calidad del corpus Gold & \texttt{aprobado: true}\\
fna-api & Endpoints con contratos y pruebas & \texttt{test\_api.py}\\
fna-pruebas & Ejecutar e interpretar pruebas y evaluación & 0 fallidas; cobertura $\geq$ 85\,\%\\
fna-seguridad & Secretos, inyección, permisos, MCP, skills de terceros & Lista de verificación con evidencia\\
fna-documentacion & README, manual, ADR, redacción impersonal & PDF compila sin referencias rotas\\
fna-simulador-3d & Piezas, articulaciones, despiece, GLB y visor (v2.0) & Pruebas del simulador y \texttt{npm test}\\
fna-maquinas-anatomia & Máquinas interactivas con datos del scraping (v2.0) & Vitest, Playwright y recorrido $<$ 30\,s\\\bottomrule
\end{longtable}
Cada skill define responsabilidad, entradas, procedimiento, salidas, restricciones y criterio de validación.
Antes de instalar skills externas: \cmd{skillspector scan <ruta-o-repositorio>}.

% ============================================================ Capítulo 15
\chapter{Gestión del proyecto con Git y GitHub}\label{cap:git}
\section{Creación del repositorio}
El proyecto se publica en \url{https://github.com/JulianReneMunoz/fabrica-nacional-autopartes} desde el
computador con Windows~11. El script \ruta{scripts/publicar_github.ps1} verifica Git y GitHub CLI, inicia
sesión con \cmd{gh auth login} si hace falta, comprueba que la cuenta activa sea \texttt{JulianReneMunoz},
inicializa la rama \texttt{main}, confirma que \texttt{.env} esté excluido, crea el primer \emph{commit} y
ejecuta \cmd{gh repo create} (público por defecto; \texttt{-Privado} para privado).
\begin{lstlisting}[style=ps]
cd C:\fabrica-nacional-autopartes
gh auth login --hostname github.com --git-protocol https --web
.\scripts\publicar_github.ps1                       # crea y publica el repositorio
.\scripts\ejecutar_pipeline.ps1 -MaxPaginas 60 -Publicar   # scraping real y evidencias en un commit propio
\end{lstlisting}
El scraping real queda en un \emph{commit} separado (\texttt{datos: scraping real \dots}), de modo que el
historial del repositorio muestra cuándo y desde dónde se obtuvieron los datos. Equivalente manual:
\begin{lstlisting}[style=ps]
git init -b main ; git add -A ; git commit -m "feat: Fábrica Nacional de Autopartes v2.0"
gh repo create JulianReneMunoz/fabrica-nacional-autopartes --public --source . --remote origin --push
\end{lstlisting}

\section{Flujo de trabajo}
Rama \texttt{main} protegida; ramas de trabajo \texttt{feat/<tema>}, \texttt{fix/<tema>} y \texttt{docs/<tema>};
\emph{commits} con prefijos convencionales (\texttt{feat:}, \texttt{fix:}, \texttt{test:}, \texttt{docs:},
\texttt{chore:}); \emph{pull request} con la CI en verde y revisión de al menos un integrante.
\begin{lstlisting}[style=ps]
git switch -c feat/scraping-real
uv run pytest -q ; uv run ruff check .
git add -A ; git commit -m "feat: captura del sitio oficial en Bronze"
git push -u origin feat/scraping-real
gh pr create --fill
uv run pre-commit install           # ganchos locales: formato, YAML/JSON, claves privadas, capas
gga init ; gga install              # opcional (Git Bash/WSL): revisión con IA contra AGENTS.md
\end{lstlisting}
\section{Integración continua}
\ruta{.github/workflows/ci.yml} ejecuta en Ubuntu y Windows: \cmd{uv sync --locked}, Ruff, mypy, verificación de
capas, construcción y validación de la base semilla, y pytest con cobertura mínima de 85\,\%; un segundo trabajo
(\texttt{visor-web}) ejecuta \cmd{npm ci}, TypeScript, Vitest y la compilación de Vite. La CI no requiere secretos
porque usa el proveedor \texttt{extractivo}.
